\section{B3: finite-grid tightness and the full contact quadratic form}
\label{sec:r33-b3}

A microscopic remainder cannot be summed over infinitely many dyadic levels. The right replacement is a finite-grid lemma with a separately proved within-cell oscillation estimate. The lemma is independent of root resampling, and hence does not alter a stopping rule. Its application to deterministic hard spheres still requires the stated moment and within-cell inputs.

\subsection{The exact Hilbert--Schmidt gap}
Let $\lambda_i\ge1$ index a mode basis and assume $\#\{i:\lambda_i\le R\}\le CR^D$. Define $H_m$ by squared coefficient weights $\lambda_i^{2m}$. The singular values of $H_{m+s}\hookrightarrow H_m$ are $\lambda_i^{-s}$, so
\[
 \sum_i\lambda_i^{-2s}
 \le C\sum_{j\ge0}2^{jD}2^{-2sj}<\infty\quad\hbox{when }2s>D.
\]
Increasing $m$ does not alter this exponent. For a contact mode space of large effective dimension, a fixed gap of two is not justified.

For $p\ge2$, if scalar increments satisfy $E|Z(e_i)|^p\le C a^p$ for all tests of unit $H_m$ norm, Minkowski in $L^{p/2}$ gives
\[
 (E\|Z\|_{H_{-(m+s)}}^p)^{2/p}
 \le\sum_i\lambda_i^{-2s}(E|Z(e_i)|^p)^{2/p}
 \le C a^2\sum_i\lambda_i^{-2s}.
\]
First apply Minkowski to finite sums and then use monotone convergence. The condition $p\ge2$ is explicit because the triangle inequality used here requires $p/2\ge1$. This supplies the actual summation argument and keeps the same final path space for all terms.

\subsection{A terminal-grid chaining theorem}
Rescale time to $[0,1]$. Let $X_\varepsilon$ be cadlag paths in a separable Hilbert space $H$. Let $J_\varepsilon\to\infty$ and $d_\varepsilon=2^{-J_\varepsilon}$. Assume for adjacent dyadic times at every level $0\le j\le J_\varepsilon$,
\begin{equation}\label{eq:r33-b3-moment}
 E\|X_\varepsilon(t+2^{-j})-X_\varepsilon(t)\|^p
 \le C[2^{-j(1+\beta)}+r_\varepsilon^p],
\end{equation}
where $p>2$, $\beta>0$. Choose $0<\gamma<\beta/p$ and require
\begin{equation}\label{eq:r33-b3-budget}
 r_\varepsilon^p d_\varepsilon^{-(1+p\gamma)}\longrightarrow0.
\end{equation}
The independent within-cell input is
\begin{equation}\label{eq:r33-b3-cell}
 \Omega_\varepsilon=
 \max_{0\le k<2^{J_\varepsilon}}
 \sup_{t\in[kd_\varepsilon,(k+1)d_\varepsilon]}
 \|X_\varepsilon(t)-X_\varepsilon(kd_\varepsilon)\|
 \longrightarrow0\quad\hbox{in probability}.
\end{equation}
Let $\widetilde X_\varepsilon$ be the linear interpolation of the terminal grid.

\begin{theorem}[Finite-grid modulus with every remainder counted]
\label{thm:r33-b3-grid}
For every $a>0$ and fixed dyadic level $J\le J_\varepsilon$,
\[
 P\{\omega_{\widetilde X_\varepsilon}(2^{-J})>C_\gamma a\}
 \le C a^{-p}\left[
 2^{-J\beta}+r_\varepsilon^p 2^{-p\gamma J}
                         2^{J_\varepsilon(1+p\gamma)}\right].
\]
Consequently $\lim_{\delta\downarrow0}\limsup_\varepsilon
P\{\omega_{X_\varepsilon}(\delta)>a\}=0$.
If the one-time laws are tight in $H$, the processes are C-tight in $D([0,1];H)$, and the interpolated laws are tight in $C([0,1];H)$.
\end{theorem}

\begin{proof}
At level $j\ge J$ set the threshold $a_j=a2^{-\gamma(j-J)}$. A union bound over its $2^j$ adjacent increments and \eqref{eq:r33-b3-moment} gives
\[
 C a^{-p}2^{-p\gamma J}
 [2^{-j(\beta-p\gamma)}+r_\varepsilon^p2^{j(1+p\gamma)}].
\]
Sum only from $J$ to $J_\varepsilon$. The first geometric sum is bounded by $Ca^{-p}2^{-J\beta}$ and the second by the displayed terminal term. On the complement of these events, any two terminal-grid times within $2^{-J}$ are joined through their dyadic ancestors; the sum of edge bounds is at most a fixed multiple of $\sum_{j\ge J}a_j=C_\gamma a$. Linear interpolation preserves this bound, up to the same fixed constant. Thus the asserted probability estimate holds.

By \eqref{eq:r33-b3-cell}, $\sup_t\|X_\varepsilon(t)-\widetilde X_\varepsilon(t)\|\le2\Omega_\varepsilon$. Use \eqref{eq:r33-b3-budget}, then let $J\to\infty$. Tightness at fixed times and this equicontinuity give path tightness by the metric Arzela--Ascoli argument. Finite-$\varepsilon$ jump processes are not called elements of $C$; it is their continuous limits and interpolants that live there.
\end{proof}

For example $p=8$, $\beta=3$, $\gamma=1/8$, $r_\varepsilon=\mu_\varepsilon^{-1/2}$ and a dyadic terminal mesh comparable to $\mu_\varepsilon^{-1}$ give a residual bound $O(\mu_\varepsilon^{-2})$. These exponents are compatible and explicitly testable.

\begin{corollary}[Stopping times are controlled without resampling]
For every temporal stopping time $\tau_\varepsilon$ and $0\le h\le\delta$,
\[
 \|X_\varepsilon((\tau_\varepsilon+h)\wedge1)-X_\varepsilon(\tau_\varepsilon)\|
 \le\omega_{X_\varepsilon}(\delta).
\]
Hence the preceding hypotheses imply the probability version of Aldous' criterion.
\end{corollary}

A small maximum jump is not a proof of \eqref{eq:r33-b3-cell}: a continuous triangular spike of height one inside a single terminal cell has zero maximum jump and vanishes at every grid point. For a decomposition with a drift $A_t=\int_0^t b_sds$, a sufficient separate drift input is $\sup_sE\|b_s\|_H^p<\infty$, which gives $E\|A_t-A_s\|_H^p\le C|t-s|^p$. Free transport must be included in this input; it is not counted as a first collision vertex.

\subsection{An unconditional static contact result}
Let $A$ be a finite positive reference contact measure and let $h\in L^2(A)$. For $\ell(q)=q\log q-q+1$, define
\[
 F_\eta(h)=\int\eta^{-2}\ell(1+\eta h)dA
\]
when $1+\eta h\ge0$, and $+\infty$ otherwise.

\begin{theorem}[Static entropy Mosco limit]
\label{thm:r33-b3-static}
As $\eta\downarrow0$, $F_\eta$ Mosco-converges on $L^2(A)$ to $\frac12\|h\|_2^2$. If a closed collision map $\mathcal C$ has closed kernel $\mathcal K$, then
\[
 \frac12\|h\|_2^2=\frac12\|h_{\mathcal K^\perp}\|_2^2+
                  \frac12\|h_{\mathcal K}\|_2^2.
\]
Only a subsequent infimum over $h$ with $\mathcal Ch=r$ removes the cycle term.
\end{theorem}

\begin{proof}
For weak $h_\eta\rightharpoonup h$ use the entropy duality with any bounded $k$:
\[
 F_\eta(h_\eta)\ge\int kh_\eta\,dA-
              \int\frac{e^{\eta k}-1-\eta k}{\eta^2}dA.
\]
The limit is $\int kh\,dA-\frac12\int k^2dA$. Taking the supremum over bounded $k$, dense in $L^2(A)$, proves the liminf. For recovery truncate $h$ at level $M$, apply the uniform Taylor expansion on that bounded set, and let $M\to\infty$ along a diagonal with $\eta M\to0$. Orthogonal decomposition proves the last formula. If the affine constraint has a solution, minimizing the norm removes its kernel component by Pythagoras.
\end{proof}

\paragraph{Mechanical application still to be established.}
A moving reference $A_{f_\eta}$, exact nonlinear balance recovery, the kinetic graph inverse and the joint microscopic CLT need further estimates. In particular the hard-sphere within-cell oscillation bound cannot be replaced by the abstract grid lemma. The static theorem is proved here; a constrained dynamic Mosco theorem or process covariance identification is not inferred merely from its formula.
